% ------------------------------------------------------------------------------
% 6. ALCANCES Y LIMITACIONES
% ------------------------------------------------------------------------------
\subsection{\pfxperfil 6. Alcances y limitaciones}\label{sec:alcances_limitaciones}

El presente proyecto abarca el diseño y desarrollo de un sistema de información web y móvil para la automatización, gestión y precalificación del proceso de postulación a la beca comedor universitaria, así como el seguimiento del beneficio otorgado. A continuación, se detallan las funciones cubiertas, los actores involucrados y las restricciones del software.

\subsubsection{\pfxperfil 6.1. Funciones principales}\label{sec:funciones_principales}

% Lista con incisos a), b), c)... para detallar cada módulo
\newenvironment{incisos}{\begin{enumerate}\renewcommand{\labelenumi}{\alph{enumi})}}{\end{enumerate}}

\noindent\textbf{Registro y autenticación:}
\begin{incisos}
    \item Permite a los postulantes, al personal administrativo de la DUBSS y a los directivos crear una cuenta e ingresar de forma segura desde el portal web o la aplicación móvil.
    \item Identifica el rol de cada usuario y lo dirige al módulo que le corresponde.
    \item Rechaza las credenciales inválidas y muestra un mensaje de error.
\end{incisos}

\noindent\textbf{Postulación en línea:}
\begin{incisos}
    \item Permite al estudiante completar el formulario socioeconómico (ingresos, dependientes, procedencia, carrera, entre otros).
    \item Permite adjuntar los documentos de respaldo digitalizados y enviar la solicitud.
    \item Valida que los datos y los documentos obligatorios estén completos y registra la fecha de envío.
\end{incisos}

\noindent\textbf{Preselección con inteligencia artificial:}
\begin{incisos}
    \item Evalúa cada solicitud recibida (20,000 o más por gestión) a partir de las variables socioeconómicas y académicas declaradas.
    \item Calcula un puntaje de vulnerabilidad por postulante y registra su valor y la fecha de cálculo para garantizar la trazabilidad.
    \item Genera una lista ordenada que prioriza a los estudiantes con mayor necesidad.
\end{incisos}

\noindent\textbf{Priorización de carpetas para auditoría:}
\begin{incisos}
    \item Ordena las solicitudes según el puntaje obtenido.
    \item Presenta al personal de la DUBSS, en primer lugar, los expedientes con mayor requerimiento económico, evitando la revisión al azar.
\end{incisos}

\noindent\textbf{Auditoría y adjudicación:}
\begin{incisos}
    \item Permite al personal de la DUBSS revisar los expedientes priorizados y contrastar los datos declarados con la documentación física de los finalistas.
    \item Permite aprobar, devolver o archivar una postulación.
    \item Adjudica la beca a los aproximadamente 900 beneficiarios y emite el acta de adjudicación.
\end{incisos}

\noindent\textbf{Asistencia (físico y digital):}
\begin{incisos}
    \item Registra la asistencia diaria de los becados al comedor.
    \item Genera las planillas y comprobantes impresos, con espacio para la firma y la fotocopia del carnet que exigen las auditorías externas.
    \item Sincroniza la información con el registro digital para permitir auditorías cruzadas.
\end{incisos}

\noindent\textbf{Consultas y notificaciones:}
\begin{incisos}
    \item Permite al estudiante consultar en tiempo real el estado de su postulación y su puntaje.
    \item Envía alertas (notificaciones push) sobre la convocatoria, la solicitud de documentación física y el resultado de la adjudicación.
\end{incisos}

\noindent\textbf{Reportes y tableros de Business Intelligence:}
\begin{incisos}
    \item Ofrece a las autoridades tableros interactivos con indicadores sobre postulaciones, eficiencia de la preselección y asistencia al comedor.
    \item Presenta proyecciones predictivas que apoyan la planificación y la rendición de cuentas.
\end{incisos}

\subsubsection{\pfxperfil 6.2. Usuarios o beneficiarios}\label{sec:usuarios_beneficiarios}

La Tabla~\ref{tab:usuarios_beneficiarios} resume los usuarios del sistema y la forma en que cada uno interactúa con él.

\begin{table}[H]
\caption{Usuarios y beneficiarios del sistema}
\label{tab:usuarios_beneficiarios}
\begin{threeparttable}
\begin{tabular}{>{\raggedright\arraybackslash}p{3.6cm} >{\raggedright\arraybackslash}p{5.2cm} >{\raggedright\arraybackslash}p{4.4cm}}
\toprule
Usuario o beneficiario & Descripción & Interacción con el sistema \\
\midrule
Estudiantes postulantes & Alumnos matriculados que solicitan la Beca Comedor. & Registran su postulación y cargan sus documentos desde la aplicación móvil o la plataforma web. \\
Estudiantes adjudicados (becados) & Beneficiarios de la beca. & Consultan el estado de su beca y gestionan sus requisitos de asistencia diaria. \\
Personal administrativo de Bienestar Estudiantil & Encargados de la gestión del proceso dentro de la DUBSS. & Revisan y auditan las carpetas priorizadas por el sistema y gestionan las convocatorias. \\
Auditores institucionales & Personal fiscalizador interno o externo. & Validan la entrega del beneficio cotejando el registro digital con los respaldos físicos (fotocopias y firmas). \\
Personal del comedor / empresa concesionaria & Encargados del servicio de alimentación. & Verifican la identidad del estudiante y participan en el registro de la asistencia diaria. \\
\bottomrule
\end{tabular}
\begin{tablenotes}
\item \textit{Nota.} Elaboración propia.
\end{tablenotes}
\end{threeparttable}
\end{table}

\subsubsection{\pfxperfil 6.3. Limitaciones y Restricciones}\label{sec:limitaciones_restricciones}

\begin{itemize}
    \item \textbf{Requisito de Auditoría Física (Restricción Normativa):} El sistema no elimina por completo el uso de papel en la fase de asistencia diaria, debido a la exigencia reglamentaria de los organismos auditores que requieren carpetas físicas con fotocopia de carnet y firma manuscrita para validar la entrega del alimento a cada becado. El software actuará como un medio de control, validación y archivo digitalizador complementario.

    \item \textbf{Operación Interna de la Cocina:} El sistema no gestiona el control de inventarios de insumos, compras de materia prima ni la contabilidad interna de la empresa privada concesionaria del comedor.

    \item \textbf{Comercialización a No Becados:} La venta de almuerzos y platos a estudiantes regulares no becados y personal universitario queda fuera del alcance del sistema, correspondiendo a la gestión propia de la empresa privada.

    \item \textbf{Dependencia de Conectividad:} La carga de expedientes y la sincronización de asistencias requieren conexión a Internet activa.

    \item \textbf{Límites de Verificación Automática:} El sistema realiza una precalificación algorítmica basada en la información declarada y digitalizada; sin embargo, la verificación presencial o visitas domiciliarias mantendrán el criterio operativo del personal de Trabajo Social.
\end{itemize}

